\documentclass[UTF8]{ctexart}
\usepackage[a4paper,margin=2.4cm]{geometry}
\usepackage{amsmath,amssymb}
\usepackage{hyperref}
\usepackage{enumitem}
\usepackage{fancyvrb}
\usepackage{xcolor}
\hypersetup{colorlinks=true,linkcolor=blue,urlcolor=blue}
\setlist{nosep,leftmargin=2em}
\DefineVerbatimEnvironment{codeblock}{Verbatim}{fontsize=\small,frame=single,rulecolor=\color{black!20}}

\title{萤火虫同步模拟：规则总表与 LaTeX 公式}
\author{}
\date{}

\begin{document}
\maketitle
\tableofcontents
\newpage

本文是当前 \texttt{\detokenize{firefly_sim}} 项目的规则说明，覆盖相位同步、移动、蝙蝠捕食者、光污染、障碍物、多种群、交互工具、性能与调试规则。它对应当前实现，而不是早期草案。

\bigskip
\hrule
\bigskip

\section{核心目标}

模拟二维空间中的萤火虫同步现象：

\[
\text{local interaction}+\text{phase oscillator}
\longrightarrow
\text{collective synchronization}
\]

每只萤火虫是一个相位振荡器。它只能和视觉范围内、且未被障碍物遮挡的邻居耦合。系统可出现：

\begin{itemize}
\item 全局无序：$r \approx 0$ 且 $\bar r_{\text{local}}\approx 0$。
\item 局部同步岛：$r \approx 0$ 但 $\bar r_{\text{local}}>0$。
\item 全局同步：$r \approx 1$ 且 $\bar r_{\text{local}}\approx 1$。
\end{itemize}

\bigskip
\hrule
\bigskip

\section{当前默认参数}

当前 UI 默认值以降低 CPU/内存压力为主：

\begin{codeblock}
N = 125
L = 10
K = 2
R_visual = 2
D = 0.02
omega0 = 1
sigma_omega = 0.5
dt = 0.01
speed = 1
sigma_flash = 0.25
mobilityEnabled = true
moveProbability = 0.1
v_firefly = 0.25
batCount = 0
\end{codeblock}

重要性能规则：

\begin{itemize}
\item Canvas 默认缩小，避免大画布造成 GPU/CPU 压力。
\item UI snapshot 约 10 fps 刷新；仿真可继续按帧推进。
\item 默认 \texttt{\detokenize{N=125}}，如需更高密度可手动调高。
\item \texttt{\detokenize{R_visual}} 可调范围比默认更大，用于补偿低 \texttt{\detokenize{N}} 时邻居过少的问题。
\end{itemize}

\bigskip
\hrule
\bigskip

\section{状态变量}

\subsection{萤火虫}

\begin{codeblock}
struct Firefly {
  float x, y;
  float vx, vy;
  float heading;
  float speed;
  float theta;
  float omega;
  float brightness;
  float localOrder;
  float panic;
  int neighborCount;
  uint8_t species;
  uint8_t alive;
};
\end{codeblock}

解释：

\begin{itemize}
\item \texttt{\detokenize{x,y}}：空间位置。
\item \texttt{\detokenize{vx,vy}}：当前速度。
\item \texttt{\detokenize{heading}}：移动方向。
\item \texttt{\detokenize{theta}}：相位。
\item \texttt{\detokenize{omega}}：自然频率。
\item \texttt{\detokenize{brightness}}：由相位或捕食者压制规则决定。
\item \texttt{\detokenize{panic}}：受蝙蝠压力影响的恐慌指标。
\item \texttt{\detokenize{alive=0}}：被蝙蝠捕获，不再参与同步和渲染。
\end{itemize}

\subsection{蝙蝠}

\begin{codeblock}
struct Bat {
  float x, y;
  float vx, vy;
  float heading;
  float speed;
  float perceptionRadius;
  float captureRadius;
  int targetIndex;
  float hunger;
  float speedBias;
  float turnRate;
  float decisionTimer;
  float decisionInterval;
  float brightnessWeight;
  float distanceWeight;
  float noisePhase;
  uint8_t active;
};
\end{codeblock}

每只蝙蝠有个体差异，避免两个蝙蝠重叠后行为完全同步。

\bigskip
\hrule
\bigskip

\section{初始条件}

萤火虫位置：

\[
x_i,y_i\sim U(0,L)
\]

相位：

\[
\theta_i(0)\sim U(0,2\pi)
\]

单种群自然频率：

\[
\omega_i\sim \mathcal{N}(\omega_0,\sigma_\omega^2)
\]

两种群自然频率：

\[
\omega_i^{(A)}\sim \mathcal{N}(\omega_A,\sigma_\omega^2),
\qquad
\omega_i^{(B)}\sim \mathcal{N}(\omega_B,\sigma_\omega^2)
\]

\bigskip
\hrule
\bigskip

\section{空间邻居与遮挡}

距离：

\[
d_{ij}=\|\mathbf{x}_i-\mathbf{x}_j\|
\]

基础邻接：

\[
A_{ij}=
\begin{cases}
1,&0<d_{ij}\le R_{\text{visual}}\\
0,&\text{otherwise}
\end{cases}
\]

邻居数量：

\[
k_i=\sum_j A_{ij}
\]

如果启用障碍物遮挡：

\[
A_{ij}
=
\mathbb{I}(0<d_{ij}\le R_{\text{visual}})
\cdot
\mathbb{I}(\text{line}(i,j)\ \text{not blocked})
\]

被捕获的萤火虫 \texttt{\detokenize{alive=0}}，不参与邻接、局部同步、全局同步和渲染。

\bigskip
\hrule
\bigskip

\section{Spatial Kuramoto 相位动力学}

单种群基础模型：

\[
\frac{d\theta_i}{dt}
=
\omega_i
+
\frac{K}{k_i}
\sum_j A_{ij}\sin(\theta_j-\theta_i)
+
\xi_i(t)
\]

若 $k_i=0$，耦合项为 $0$。归一化因子 $1/k_i$ 防止邻居多的个体受到过强耦合。

噪声：

\[
\xi_i(t)=\sqrt{2D}\eta_i(t)
\]

随机微分形式：

\[
d\theta_i
=
\left[
\omega_i+
\frac{K}{k_i}
\sum_j A_{ij}\sin(\theta_j-\theta_i)
\right]dt
+
\sqrt{2D}\,dW_i(t)
\]

\bigskip
\hrule
\bigskip

\section{数值更新}

Euler-Maruyama 离散化：

\[
\theta_i(t+\Delta t)
=
\theta_i(t)
+
\left[
\omega_i+
\frac{K}{k_i}
\sum_j A_{ij}\sin(\theta_j(t)-\theta_i(t))
\right]\Delta t
+
\sqrt{2D\Delta t}Z_i
\]

其中：

\[
Z_i\sim\mathcal{N}(0,1)
\]

相位 wrap：

\[
\theta_i\leftarrow \theta_i\bmod 2\pi
\]

\bigskip
\hrule
\bigskip

\section{闪光与亮度}

二值闪光：

\[
F_i(t)=
\begin{cases}
1,&\theta_i(t-\Delta t)>\theta_i(t)\\
0,&\text{otherwise}
\end{cases}
\]

余弦亮度：

\[
B_i(t)=\frac{1+\cos\theta_i(t)}{2}
\]

尖峰亮度：

\[
B_i(t)=
\exp\left[
-\frac{\operatorname{wrap}(\theta_i(t))^2}
{2\sigma_{\text{flash}}^2}
\right]
\]

其中：

\[
\operatorname{wrap}(\theta)=((\theta+\pi)\bmod 2\pi)-\pi
\]

\subsection{蝙蝠压制闪光规则}

如果活萤火虫处在任意蝙蝠感知半径内：

\[
d_{ib}\le R_{\text{bat\_perception}}
\]

则本步：

\[
\theta_i(t+\Delta t)=\theta_i(t),
\qquad
B_i(t+\Delta t)=0
\]

也就是说：它不发光，也不更新自然频率/相位。

\bigskip
\hrule
\bigskip

\section{同步指标}

全局 Kuramoto order parameter：

\[
r(t)e^{i\psi(t)}
=
\frac{1}{N_{\text{alive}}}
\sum_{j:\text{alive}} e^{i\theta_j(t)}
\]

\[
r(t)
=
\left|
\frac{1}{N_{\text{alive}}}
\sum_{j:\text{alive}} e^{i\theta_j(t)}
\right|
\]

局部同步：

\[
r_i^{\text{local}}(t)
=
\left|
\frac{1}{k_i}
\sum_j A_{ij}e^{i\theta_j(t)}
\right|
\]

平均局部同步：

\[
\bar r_{\text{local}}(t)
=
\frac{1}{N_{\text{alive}}}
\sum_{i:\text{alive}} r_i^{\text{local}}(t)
\]

\bigskip
\hrule
\bigskip

\section{临界耦合与扫描}

稳态平均：

\[
\bar r(q)=
\frac{1}{M_T}
\sum_{t>T_{\text{burn}}}
r(t)
\]

临界耦合估计：

\[
K_c=\min\{K:\bar r(K)>r_{\text{threshold}}\}
\]

全局耦合高斯频率分布的参考近似：

\[
K_c=
\frac{2\sqrt{2\pi}\sigma_\omega}{\pi}
\]

空间局部耦合下实际临界值依赖：

\[
K_c=K_c(\sigma_\omega,R_{\text{visual}},N,D)
\]

\bigskip
\hrule
\bigskip

\section{城市光污染}

城市灯光作为外部驱动：

\[
\frac{d\theta_i}{dt}
=
\omega_i
+
\frac{K}{k_i}
\sum_j A_{ij}\sin(\theta_j-\theta_i)
+
\epsilon_{\text{city}}
\sin(\Omega_{\text{city}}t+\phi_{\text{city}}-\theta_i)
+
\xi_i(t)
\]

锁频指标：

\[
\Delta_{\text{lock}}
=
\left|
\frac{d\psi}{dt}
-
\Omega_{\text{city}}
\right|
\]

\bigskip
\hrule
\bigskip

\section{多种群规则}

种内/种间耦合：

\[
K_{ij}
=
\begin{cases}
K_{\text{in}},&s_i=s_j\\
K_{\text{out}},&s_i\ne s_j
\end{cases}
\]

两种群动力学：

\[
\frac{d\theta_i}{dt}
=
\omega_i+
\frac{1}{k_i}
\sum_j A_{ij}K_{ij}\sin(\theta_j-\theta_i)
+
\xi_i(t)
\]

种群序参量：

\[
r_A(t)=
\left|
\frac{1}{N_A}
\sum_{j\in A}e^{i\theta_j(t)}
\right|,
\qquad
r_B(t)=
\left|
\frac{1}{N_B}
\sum_{j\in B}e^{i\theta_j(t)}
\right|
\]

\bigskip
\hrule
\bigskip

\section{萤火虫移动规则}

萤火虫移动是概率事件，不是每步都移动。

普通状态：

\[
P(\text{move at step})=p_{\text{move}}
\]

当前默认：

\[
p_{\text{move}}=0.1
\]

如果被蝙蝠感知/追逐：

\[
d_{ib}\le R_{\text{bat\_perception}}
\Longrightarrow
P(\text{move at step})=1
\]

移动方向使用 correlated random walk：

\[
\alpha_i(t+\Delta t)
=
\alpha_i(t)
+
\sqrt{2D_{\text{turn}}\Delta t}Z_i
\]

随机移动速度：

\[
\mathbf{v}_i^{\text{random}}
=
v_{\text{firefly}}
\begin{bmatrix}
\cos\alpha_i\\
\sin\alpha_i
\end{bmatrix}
+
\sqrt{2D_{\text{move}}\Delta t}\boldsymbol{\eta}_i
\]

\bigskip
\hrule
\bigskip

\section{蝙蝠躲避规则}

对满足 $d_{ib}<R_{\text{avoid}}$ 的蝙蝠，萤火虫加入排斥速度：

\[
\mathbf{v}_i^{\text{avoid}}
=
\sum_b
\chi_{\text{bat}}
\left(1-\frac{d_{ib}}{R_{\text{avoid}}}\right)
\frac{\mathbf{x}_i-\mathbf{x}_b}{d_{ib}+\epsilon}
\]

总速度：

\[
\mathbf{v}_i
=
\mathbf{v}_i^{\text{random}}
+
\mathbf{v}_i^{\text{avoid}}
\]

位置更新：

\[
\mathbf{x}_i(t+\Delta t)
=
\mathbf{x}_i(t)
+
\mathbf{v}_i\Delta t
\]

边界规则为反射边界，保证所有 agent 留在 $[0,L]^2$。

恐慌指标：

\[
p_i
=
\max_b
\left(1-\frac{d_{ib}}{R_{\text{avoid}}}\right)_+
\]

\bigskip
\hrule
\bigskip

\section{蝙蝠捕食者策略}

\subsection{个体差异}

每只蝙蝠初始化时带有固定个体参数：

\begin{codeblock}
speedBias
turnRate
decisionInterval
brightnessWeight
distanceWeight
noisePhase
\end{codeblock}

实际速度：

\[
v_b^{\text{actual}}
=
v_{\text{bat}}\cdot \text{speedBias}_b
\]

\subsection{top-k softmax 目标选择}

蝙蝠不是每步都重新决策，而是每隔自己的 \texttt{\detokenize{decisionInterval}} 重新选择目标。

候选萤火虫需满足：

\[
d_{bj}\le R_{\text{bat\_perception}},
\qquad
\text{alive}_j=1
\]

目标得分：

\[
s_{bj}
=
w_b^{(B)}
\frac{B_j}{d_{bj}+\epsilon}
+
w_b^{(D)}
\frac{1}{d_{bj}+\epsilon}
+
\zeta
\]

其中：

\begin{itemize}
\item $w_b^{(B)}$ 对应 \texttt{\detokenize{brightnessWeight}}。
\item $w_b^{(D)}$ 对应 \texttt{\detokenize{distanceWeight}}。
\item $\zeta$ 是小随机扰动。
\end{itemize}

先取分数最高的 top-k：

\[
\mathcal{C}_b=\operatorname{TopK}_j(s_{bj})
\]

再用 softmax 采样：

\[
P(j\mid b)
=
\frac{
\exp(s_{bj}/T_{\text{bat}})
}{
\sum_{\ell\in \mathcal{C}_b}
\exp(s_{b\ell}/T_{\text{bat}})
}
\]

\texttt{\detokenize{T_bat = batSoftmaxTemperature}}。温度越低越接近贪心 argmax，温度越高越随机。

\subsection{追逐更新}

目标方向：

\[
\beta_b^*
=
\operatorname{atan2}(y_{j^*}-y_b,x_{j^*}-x_b)
+
\sigma_{\text{chase}}Z_b
\]

平滑转向：

\[
\beta_b(t+\Delta t)
=
\beta_b(t)
+
\lambda_b
\operatorname{wrap}(\beta_b^*-\beta_b(t))
\]

其中 $\lambda_b$ 对应 \texttt{\detokenize{turnRate}}。

\subsection{蝙蝠-蝙蝠分离}

当两个蝙蝠距离过近：

\[
d_{bb'}<R_{\text{bat\_sep}}
\]

加入分离速度：

\[
\mathbf{v}_b^{\text{sep}}
=
\sum_{b'\ne b}
\gamma_{\text{sep}}
\left(1-\frac{d_{bb'}}{R_{\text{bat\_sep}}}\right)
\frac{\mathbf{x}_b-\mathbf{x}_{b'}}{d_{bb'}+\epsilon}
\]

这样即使两个蝙蝠重叠，也会因为个体参数、softmax 采样、决策间隔和分离力而快速分化。

\subsection{捕获规则}

如果：

\[
d_{ib}<R_{\text{capture}}
\]

则：

\[
\text{alive}_i=0
\]

被捕获的萤火虫不再参与：

\begin{itemize}
\item 相位更新
\item 亮度显示
\item 邻居图
\item order parameter
\item local order
\end{itemize}

\bigskip
\hrule
\bigskip

\section{交互工具规则}

工具栏：

\begin{itemize}
\item \texttt{\detokenize{+}}：添加萤火虫。
\item \texttt{\detokenize{-}}：通用橡皮擦，擦除半径内的萤火虫、障碍物、城市光源和蝙蝠。
\item \texttt{\detokenize{O}}：添加障碍物。
\item \texttt{\detokenize{*}}：添加城市光源。
\item \texttt{\detokenize{B}}：添加蝙蝠。
\item \texttt{\detokenize{Inspect}}：悬停查看单个萤火虫状态。
\end{itemize}

数量联动：

\begin{itemize}
\item \texttt{\detokenize{N}} slider 改变会同步实际萤火虫数量。
\item 画布添加/擦除萤火虫也会同步回 \texttt{\detokenize{N}} slider。
\item 添加/清除/擦除蝙蝠会同步回 \texttt{\detokenize{batCount}}。
\end{itemize}

\bigskip
\hrule
\bigskip

\section{当前 metrics}

核心 metrics：

\begin{codeblock}
r
psi
rLocalMean
avgNeighbors
isolatedCount
flashCount
cityLockDelta
rA
rB
aliveCount
capturedCount
meanPanic
meanNearestBatDistance
batTargetCount
\end{codeblock}

解释：

\begin{itemize}
\item \texttt{\detokenize{aliveCount}}：仍存活的萤火虫数量。
\item \texttt{\detokenize{capturedCount}}：已被捕获的萤火虫数量。
\item \texttt{\detokenize{meanPanic}}：平均恐慌程度。
\item \texttt{\detokenize{meanNearestBatDistance}}：活萤火虫到最近蝙蝠的平均距离。
\item \texttt{\detokenize{batTargetCount}}：当前有目标的蝙蝠数量。
\end{itemize}

\bigskip
\hrule
\bigskip

\section{性能与调试规则}

为避免浏览器 out-of-memory：

\begin{itemize}
\item 默认 \texttt{\detokenize{N=125}}。
\item 默认 \texttt{\detokenize{speed=1}}。
\item Canvas 尺寸受限。
\item UI snapshot 约 10 fps 更新；暂停时约 2 fps。
\item diagnostics 最多保留 80 条。
\item info 级日志不写入 console，避免 DevTools 保留大对象。
\end{itemize}

崩溃诊断：

\begin{itemize}
\item 捕获 browser error。
\item 捕获 unhandled promise rejection。
\item 捕获 React render crash。
\item 主循环 \texttt{\detokenize{step()}} 或 \texttt{\detokenize{getSnapshot()}} 抛错时自动暂停。
\item 右下角 \texttt{\detokenize{Logs}} 面板显示最近事件和错误上下文。
\end{itemize}

\bigskip
\hrule
\bigskip

\section{当前实验流程}

\subsection{Task 1：基础同步}

\begin{enumerate}
\item 设置低或中等 \texttt{\detokenize{N}}。
\item 调整 \texttt{\detokenize{K}} 与 \texttt{\detokenize{R_visual}}。
\item 观察 $r(t)$ 和 $\bar r_{\text{local}}(t)$。
\end{enumerate}

\subsection{Task 2：寻找临界点}

扫描 $K$：

\[
\bar r(K)
=
\frac{1}{M_T}
\sum_{t>T_{\text{burn}}}
r(t)
\]

估计：

\[
K_c=\min\{K:\bar r(K)>r_{\text{threshold}}\}
\]

\subsection{Task 3：光污染}

调整：

\[
\epsilon_{\text{city}},
\qquad
\Omega_{\text{city}}
\]

观察 \texttt{\detokenize{Delta_lock}}。

\subsection{Task 4：树林遮挡}

添加障碍物并观察：

\[
r(t)
\quad\text{vs.}\quad
\bar r_{\text{local}}(t)
\]

\subsection{Task 5：蝙蝠压力}

添加蝙蝠，观察：

\begin{codeblock}
aliveCount
capturedCount
meanPanic
meanNearestBatDistance
batTargetCount
r(t)
\end{codeblock}

重点现象：

\begin{itemize}
\item 蝙蝠感知范围内萤火虫会熄灭并冻结相位。
\item 被追逐萤火虫一定移动。
\item 多个蝙蝠因为 softmax 策略和 separation 不会长期完全重叠。
\end{itemize}

\bigskip
\hrule
\bigskip

\section{一页 slide 公式}

主模型：

\[
\frac{d\theta_i}{dt}
=
\omega_i
+
\frac{K}{k_i}
\sum_j A_{ij}\sin(\theta_j-\theta_i)
+
\xi_i(t)
\]

序参量：

\[
r(t)=
\left|
\frac{1}{N_{\text{alive}}}
\sum_{j:\text{alive}}
e^{i\theta_j(t)}
\right|
\]

移动概率：

\[
P(\text{move})=
\begin{cases}
1,&d_{ib}\le R_{\text{bat\_perception}}\\
p_{\text{move}},&\text{otherwise}
\end{cases}
\]

蝙蝠 softmax 目标选择：

\[
P(j\mid b)
=
\frac{
\exp(s_{bj}/T_{\text{bat}})
}{
\sum_{\ell\in \mathcal{C}_b}
\exp(s_{b\ell}/T_{\text{bat}})
}
\]

一句总结：

\[
\text{synchronization}
\quad\text{emerges from}\quad
\text{local coupling},
\quad
\text{but predators reshape who flashes and who moves.}
\]

\end{document}
